% GENERATED FILE. Edit canonical lessons, then run npm run generate:books.
% canonical-source-hash: fnv1a64:36967e369abf8ead
% canonical-lessons: SA-C414-digbhramah, SA-C414-bhrantah, SA-C414-pratipadyate, SA-C414-aparicitah, SA-C414-agantukah

\chapter{Losing the Way}
\label{ch:sa-losing-the-way}
\hlchaptermodality{414}

\begin{chapteropening}
\textbf{By the end of this chapter:} \emph{I can say I am lost and find the way again.}

Complete the last lesson of chapter 414: I can say I am lost and find the way again.
\end{chapteropening}

\section[\texorpdfstring{\sk{दिग्भ्रमः} (digbhramaḥ)}{दिग्भ्रमः (digbhramaḥ)}]{\sk{दिग्भ्रमः} (digbhramaḥ) — losing one's bearings}

\label{lesson:SA-C414-digbhramah}



\begin{quote}
Before the new one: say the Sanskrit for delivers, then the Sanskrit for arranges.
\end{quote}

\subsection*{You'll want to know: \sk{दिग्भ्रमः}}
\textbf{\sk{दिग्भ्रमः}} — \emph{digbhramaḥ} — \textquotedblleft{}losing one's bearings\textquotedblright{}.

\begin{quote}
\sk{वने} \sk{पथिकस्य} \sk{दिग्भ्रमः} \sk{भवति।}
\end{quote}

\begin{quote}
\emph{vane pathikasya digbhramaḥ bhavati.}
\end{quote}

\begin{quote}
In the forest the wayfarer loses his bearings.
\end{quote}

A masculine noun.

\begin{cousinweb}
\textbf{\sk{दिग्भ्रमः}} (\emph{digbhramaḥ}) is built from \emph{diś}, \textquotedblleft{}a direction\textquotedblright{}, and \emph{bhrama}, \textquotedblleft{}wandering, confusion\textquotedblright{}, from the root √bhram, \textquotedblleft{}to wander\textquotedblright{}; MW gives the compound as confusion about the quarters of the compass, mistaking the way (MW).
\end{cousinweb}

\begin{grammarlens}[title={What you've built}]
One more word for getting around.
\end{grammarlens}

\subsection*{Your turn}
\begin{itemize}
\raggedright
  \item \emph{Say it:} \emph{digbhramaḥ}
  \item \emph{Say it:} \emph{digbhramaḥ}, once more
  \item \emph{Say it:} \emph{vane pathikasya digbhramaḥ bhavati.} — In the forest the wayfarer loses his bearings.
  \item \emph{Recall:} say the Sanskrit for delivers, then the Sanskrit for arranges, then say \emph{digbhramaḥ} again
\end{itemize}

\begin{tcolorbox}[breakable,colback=teal!4,colframe=teal!35!black,title={Before you move on}]
What does \sk{दिग्भ्रमः} mean? (\textquotedblleft{}Losing one's bearings\textquotedblright{}.) Say it once more.
\end{tcolorbox}
\section[\texorpdfstring{\sk{भ्रान्तः} (bhrāntaḥ)}{भ्रान्तः (bhrāntaḥ)}]{\sk{भ्रान्तः} (bhrāntaḥ) — astray; wandering}

\label{lesson:SA-C414-bhrantah}



\begin{quote}
Before the new one: say the Sanskrit for arranges, then the Sanskrit for losing one's bearings.
\end{quote}

\subsection*{You'll want to know: \sk{भ्रान्तः}}
\textbf{\sk{भ्रान्तः}} — \emph{bhrāntaḥ} — \textquotedblleft{}astray; wandering\textquotedblright{}.

\begin{quote}
\sk{भ्रान्तः} \sk{पथिकः} \sk{ग्रामं} \sk{न} \sk{प्राप्नोति।}
\end{quote}

\begin{quote}
\emph{bhrāntaḥ pathikaḥ grāmaṃ na prāpnoti.}
\end{quote}

\begin{quote}
The wayfarer who has gone astray does not reach the village.
\end{quote}

\begin{cousinweb}
\textbf{\sk{भ्रान्तः}} (\emph{bhrāntaḥ}) comes from the root √bhram, \textquotedblleft{}to wander\textquotedblright{}, which also gives \emph{bhramati}, \textquotedblleft{}wanders\textquotedblright{}: it is the participle, \textquotedblleft{}having wandered\textquotedblright{}, and MW gives \emph{bhrānta} as roaming about and as confused (MW).
\end{cousinweb}

\begin{grammarlens}[title={What you've built}]
One more word for getting around.
\end{grammarlens}

\subsection*{Your turn}
\begin{itemize}
\raggedright
  \item \emph{Say it:} \emph{bhrāntaḥ}
  \item \emph{Say it:} \emph{bhrāntaḥ}, once more
  \item \emph{Say it:} \emph{bhrāntaḥ pathikaḥ grāmaṃ na prāpnoti.} — The wayfarer who has gone astray does not reach the village.
  \item \emph{Recall:} say the Sanskrit for arranges, then the Sanskrit for losing one's bearings, then say \emph{bhrāntaḥ} again
\end{itemize}

\begin{tcolorbox}[breakable,colback=teal!4,colframe=teal!35!black,title={Before you move on}]
What does \sk{भ्रान्तः} mean? (\textquotedblleft{}Astray; wandering\textquotedblright{}.) Say it once more.
\end{tcolorbox}
\section[\texorpdfstring{\sk{प्रतिपद्यते} (pratipadyate)}{प्रतिपद्यते (pratipadyate)}]{\sk{प्रतिपद्यते} (pratipadyate) — finds the way; reaches}

\label{lesson:SA-C414-pratipadyate}



\begin{quote}
Before the new one: say the Sanskrit for losing one's bearings, then the Sanskrit for astray.
\end{quote}

\subsection*{You'll want to know: \sk{प्रतिपद्यते}}
\textbf{\sk{प्रतिपद्यते}} — \emph{pratipadyate} — \textquotedblleft{}finds the way; reaches\textquotedblright{}.

\begin{quote}
\sk{पथिकः} \sk{शनैः} \sk{मार्गं} \sk{प्रतिपद्यते।}
\end{quote}

\begin{quote}
\emph{pathikaḥ śanaiḥ mārgaṃ pratipadyate.}
\end{quote}

\begin{quote}
Slowly the wayfarer finds the road.
\end{quote}

Past (he …): \sk{सः} \sk{मार्गं} \sk{प्रतिपन्नवान्}, \emph{saḥ mārgaṃ pratipannavān}.

\begin{cousinweb}
\textbf{\sk{प्रतिपद्यते}} (\emph{pratipadyate}) is built from \emph{prati}, \textquotedblleft{}towards\textquotedblright{}, and the root √pad, \textquotedblleft{}to go, to step\textquotedblright{}; MW gives \emph{prati-pad} as to set foot upon, arrive at, reach, with the participle \emph{pratipanna}, \textquotedblleft{}arrived, found\textquotedblright{} (MW).
\end{cousinweb}

\begin{grammarlens}[title={What you've built}]
One more word for getting around.
\end{grammarlens}

\subsection*{Your turn}
\begin{itemize}
\raggedright
  \item \emph{Say it:} \emph{pratipadyate}
  \item \emph{Say it:} \emph{pratipadyate}, once more
  \item \emph{Say it:} \emph{pathikaḥ śanaiḥ mārgaṃ pratipadyate.} — Slowly the wayfarer finds the road.
  \item \emph{Recall:} say the Sanskrit for losing one's bearings, then the Sanskrit for astray, then say \emph{pratipadyate} again
\end{itemize}

\begin{tcolorbox}[breakable,colback=teal!4,colframe=teal!35!black,title={Before you move on}]
What does \sk{प्रतिपद्यते} mean? (\textquotedblleft{}Finds the way; reaches\textquotedblright{}.) Say it once more.
\end{tcolorbox}
\section[\texorpdfstring{\sk{अपरिचितः} (aparicitaḥ)}{अपरिचितः (aparicitaḥ)}]{\sk{अपरिचितः} (aparicitaḥ) — unfamiliar; strange}

\label{lesson:SA-C414-aparicitah}



\begin{quote}
Before the new one: say the Sanskrit for astray, then the Sanskrit for finds the way.
\end{quote}

\subsection*{You'll want to know: \sk{अपरिचितः}}
\textbf{\sk{अपरिचितः}} — \emph{aparicitaḥ} — \textquotedblleft{}unfamiliar; strange\textquotedblright{}.

\begin{quote}
\sk{एतत्} \sk{नगरं} \sk{मम} \sk{अपरिचितम्} \sk{अस्ति।}
\end{quote}

\begin{quote}
\emph{etat nagaraṃ mama aparicitam asti.}
\end{quote}

\begin{quote}
This city is unfamiliar to me.
\end{quote}

\begin{cousinweb}
\textbf{\sk{अपरिचितः}} (\emph{aparicitaḥ}) is built from \emph{a}, \textquotedblleft{}not\textquotedblright{}, and \emph{paricita}, \textquotedblleft{}known, familiar\textquotedblright{}, from \emph{pari} and the root √ci, \textquotedblleft{}to gather\textquotedblright{}; MW gives \emph{aparicita} as unacquainted, unknown (MW).
\end{cousinweb}

\begin{grammarlens}[title={What you've built}]
One more word for getting around.
\end{grammarlens}

\subsection*{Your turn}
\begin{itemize}
\raggedright
  \item \emph{Say it:} \emph{aparicitaḥ}
  \item \emph{Say it:} \emph{aparicitaḥ}, once more
  \item \emph{Say it:} \emph{etat nagaraṃ mama aparicitam asti.} — This city is unfamiliar to me.
  \item \emph{Recall:} say the Sanskrit for astray, then the Sanskrit for finds the way, then say \emph{aparicitaḥ} again
\end{itemize}

\begin{tcolorbox}[breakable,colback=teal!4,colframe=teal!35!black,title={Before you move on}]
What does \sk{अपरिचितः} mean? (\textquotedblleft{}Unfamiliar; strange\textquotedblright{}.) Say it once more.
\end{tcolorbox}
\section[\texorpdfstring{\sk{आगन्तुकः} (āgantukaḥ)}{आगन्तुकः (āgantukaḥ)}]{\sk{आगन्तुकः} (āgantukaḥ) — a newcomer; a stranger}

\label{lesson:SA-C414-agantukah}



\begin{quote}
Before the new one: say the Sanskrit for finds the way, then the Sanskrit for unfamiliar.
\end{quote}

\subsection*{You'll want to know: \sk{आगन्तुकः}}
\textbf{\sk{आगन्तुकः}} — \emph{āgantukaḥ} — \textquotedblleft{}a newcomer; a stranger\textquotedblright{}.

\begin{quote}
\sk{आगन्तुकः} \sk{जनान्} \sk{मार्गं} \sk{पृच्छति।}
\end{quote}

\begin{quote}
\emph{āgantukaḥ janān mārgaṃ pṛcchati.}
\end{quote}

\begin{quote}
The newcomer asks people the way.
\end{quote}

A masculine noun.

\begin{cousinweb}
\textbf{\sk{आगन्तुकः}} (\emph{āgantukaḥ}) is built from \emph{ā}, \textquotedblleft{}towards\textquotedblright{}, and the root √gam, \textquotedblleft{}to go\textquotedblright{}: one who has come; MW gives \emph{āgantuka} as a new comer, a stranger, a guest (MW).
\end{cousinweb}

\begin{grammarlens}[title={What you've built}]
One more word for getting around.
\end{grammarlens}

\subsection*{Your turn}
\begin{itemize}
\raggedright
  \item \emph{Say it:} \emph{āgantukaḥ}
  \item \emph{Say it:} \emph{āgantukaḥ}, once more
  \item \emph{Say it:} \emph{āgantukaḥ janān mārgaṃ pṛcchati.} — The newcomer asks people the way.
  \item \emph{Recall:} say the Sanskrit for finds the way, then the Sanskrit for unfamiliar, then say \emph{āgantukaḥ} again
\end{itemize}

\begin{tcolorbox}[breakable,colback=teal!4,colframe=teal!35!black,title={Before you move on}]
What does \sk{आगन्तुकः} mean? (\textquotedblleft{}A newcomer; a stranger\textquotedblright{}.) Say it once more.
\end{tcolorbox}
